\section{The RTL-to-GDS pipeline}

\begin{frame}{RTL and GDS -- two ends of the same pipeline}
  \begin{columns}[T,onlytextwidth]
    \column{0.48\textwidth}
      \textbf{RTL}\\[0.3em]
      A text description of a circuit in Verilog / SystemVerilog.\\
      Says \emph{what} the circuit does.
    \column{0.48\textwidth}
      \textbf{GDS / GDSII}\\[0.3em]
      A binary layout file -- every polygon on every layer.\\
      What the foundry turns into masks.
  \end{columns}
  \vspace{0.8em}
  \begin{center}
    RTL-to-GDS = automating the path from one to the other.
  \end{center}
  \note{%
    Audience calibration. Most first-timers understand code (RTL) and know
    chips are ``physical'' (GDS) but have never seen the in-between steps.
    Emphasise the distinction: \emph{behaviour} vs \emph{geometry}. One wrong
    polygon and the mask is unusable; one wrong environment variable and the
    tool cannot even start. The whole deck earns its keep by showing this
    translation is now reproducible in one command.\par
    Software analogy (spoken, not on the slide): RTL is like C source,
    GDS is like an ELF binary, and LibreLane is a compiler chain of $\sim$80
    passes. Unlike software, each pass can veto manufacturability.}
\end{frame}

\begin{frame}{The pipeline, at a glance}
  \begin{center}
  \begin{tikzpicture}[node distance=0.7cm and 0.35cm]
    \node[pipelinebox]  (rtl)      {RTL\\\scriptsize Verilog};
    \node[pipelinebox, right=of rtl]  (syn)   {Synthesis};
    \node[pipelinebox, right=of syn]  (fp)    {Floorplan};
    \node[pipelinebox, right=of fp]   (pdn)   {Power\\grid (PDN)};
    \node[pipelinebox, below=of pdn]  (pl)    {Placement};
    \node[pipelinebox, left=of pl]    (cts)   {Clock tree\\(CTS)};
    \node[pipelinebox, left=of cts]   (rt)    {Routing};
    \node[pipelinebox3, left=of rt]   (so)    {Signoff\\\scriptsize DRC/LVS/STA};
    \node[pipelinebox2, below=of so]  (gds)   {GDSII};
    \draw[arrow] (rtl) -- (syn);
    \draw[arrow] (syn) -- (fp);
    \draw[arrow] (fp)  -- (pdn);
    \draw[arrow] (pdn) -- (pl);
    \draw[arrow] (pl)  -- (cts);
    \draw[arrow] (cts) -- (rt);
    \draw[arrow] (rt)  -- (so);
    \draw[arrow] (so)  -- (gds);
  \end{tikzpicture}
  \end{center}
  \note{%
    Reading guide for the diagram.\par
    Synthesis (Yosys) turns RTL into a gate-level netlist of standard-cell
    instances. Floorplan decides die size and where macros / I/O sit. PDN
    draws VDD/VSS straps and rings. Placement legalises every cell to a
    standard-cell row. CTS balances the clock tree. Routing draws every
    signal wire. Signoff verifies manufacturability (DRC), schematic fidelity
    (LVS), and timing (STA).\par
    Emphasise: each stage writes a milestone artifact, and any stage can
    halt the flow. LibreLane is the thing that stitches all 80-ish steps
    together -- the rest of this deck is about driving \emph{it}, not the
    individual tools.}
\end{frame}

\begin{frame}{What you get at the end}
  \begin{columns}[T,onlytextwidth]
    \column{0.55\textwidth}
      A successful flow run hands you:
      \begin{itemize}\small
        \item \cmd{final/gds/chip\_top.gds} -- the layout.
        \item \cmd{final/def/} -- routed netlist + placement.
        \item \cmd{final/lib/} -- timing libraries (extracted).
        \item \cmd{final/spef/} -- parasitics for post-layout STA.
        \item \cmd{final/sdf/} -- back-annotation for gate-level sim.
        \item \cmd{metrics.csv} -- one row per metric, for CI diffs.
      \end{itemize}
    \column{0.45\textwidth}
      \screenshotplaceholder{slides/figures/final_layout_render.png}{%
        Final GDS rendered by KLayout.}
  \end{columns}
  \note{%
    This frame sets the "you will have something tangible" expectation.
    The final layout render is what people put on posters. Everything else
    (DEF, SPEF, SDF) matters only to the people doing follow-up work:
    gate-level sim, analog post-layout, shuttle submission.\par
    Name-drop \cmd{metrics.csv} -- that is the single file you grep to know
    whether the chip is green. Mention: commit a trimmed \cmd{metrics.csv}
    to git as a regression baseline; you will thank yourself when the image
    is rebuilt weeks later.}
\end{frame}
